% B_terminology.tex -- concise English--Vietnamese WLAN terminology
% =====================================================================
\section{Thuật ngữ chuyên ngành Anh--Việt}
\label{app:terms}

\begingroup
\small
\newcolumntype{L}[1]{>{\hsize=#1\hsize\raggedright\arraybackslash}X}
\newcommand{\catrow}[1]{%
  \multicolumn{3}{@{}l@{}}{\textbf{\large #1}} \\*
  \midrule
}

\begin{xltabular}{\linewidth}{@{} L{0.8} L{0.8} L{1.4} @{}}
  \caption{Các thuật ngữ WLAN Anh--Việt thường dùng.}
  \label{tab:terms}\\
  \toprule
  \textbf{Thuật ngữ tiếng Anh (viết tắt)} & \textbf{Tiếng Việt} &
  \textbf{Giải thích ngắn}\\
  \midrule
  \endfirsthead
  \multicolumn{3}{@{}l}{\small\textsl{Bảng~\thetable\ (tiếp theo)}}\\
  \toprule
  \textbf{Thuật ngữ tiếng Anh (viết tắt)} & \textbf{Tiếng Việt} &
  \textbf{Giải thích ngắn}\\
  \midrule
  \endhead
  \midrule
  \multicolumn{3}{r@{}}{\small\textsl{tiếp tục ở trang sau}}\\
  \endfoot
  \bottomrule
  \endlastfoot

  \catrow{Mạng và kiến trúc}
  Local Area Network (LAN) & Mạng cục bộ & Mạng truyền thông phục vụ trong phạm vi địa lý giới hạn\\
  Wireless Local Area Network (WLAN) & Mạng cục bộ không dây & Mạng cục bộ sử dụng sóng vô tuyến thay cho đường dây cáp\\
  Wi-Fi & Mạng không dây Wi-Fi & Nhãn hiệu của Wi-Fi Alliance chứng nhận cho các thiết bị đạt chuẩn 802.11\\
  Access Point (AP) & Điểm truy cập & Thiết bị bắc cầu kết nối các trạm không dây vào mạng có dây\\
  Wireless station (STA) & Trạm không dây & Thiết bị được trang bị bộ thu phát vô tuyến theo chuẩn 802.11\\
  Wireless LAN Controller (WLC) & Bộ điều khiển mạng không dây & Thiết bị quản lý và kiểm soát tập trung nhiều AP dạng lightweight\\
  Basic Service Set (BSS) & Tập dịch vụ cơ bản & Một AP cùng tập hợp các trạm không dây đã thiết lập liên kết với AP đó\\
  BSS Identifier (BSSID) & Mã định danh BSS & Mã nhận diện duy nhất cho một BSS, thông thường là địa chỉ MAC của AP\\
  Extended Service Set (ESS) & Tập dịch vụ mở rộng & Hệ thống gồm nhiều BSS liên kết qua mạng phân phối và dùng chung SSID\\
  Service Set Identifier (SSID) & Tên mạng & Chuỗi tên định danh logic được mạng không dây quảng bá\\
  Association & Liên kết & Thủ tục trạm không dây gia nhập và thiết lập liên kết với một AP\\
  Distribution System (DS) & Hệ thống phân phối & Đường trục kết nối liên thông các AP trong mạng mở rộng\\
  Roaming & Chuyển vùng & Quá trình trạm di chuyển giữa các AP mà vẫn duy trì kết nối mạng liên tục\\
  Power over Ethernet (PoE) & Cấp nguồn qua Ethernet & Công nghệ cấp nguồn điện trực tiếp cho thiết bị qua sợi cáp mạng\\
  Virtual LAN (VLAN) & Mạng LAN ảo & Kỹ thuật phân đoạn logic một mạng cục bộ vật lý thành các miền quảng bá riêng\\

  \catrow{Vô tuyến và lớp vật lý}
  Wireless medium & Môi trường truyền không dây & Môi trường truyền dẫn vô tuyến dùng chung để trao đổi khung dữ liệu\\
  Frequency & Tần số & Số chu kỳ dao động toàn phần của sóng điện từ trong một giây\\
  Band & Băng tần & Dải tần số vô tuyến được phân bổ sử dụng cho một hệ thống thông tin\\
  Channel & Kênh & Phân đoạn tần số cụ thể bên trong băng tần phục vụ một liên kết mạng\\
  Channel width & Độ rộng kênh & Bề rộng dải phổ tần số mà một kênh truyền vô tuyến chiếm dụng\\
  ISM band & Băng tần công nghiệp, khoa học và y tế & Các dải tần số không cần cấp phép như 2,4~GHz và 5~GHz\\
  Regulatory domain & Miền quản lý tần số & Khu vực quốc gia áp dụng các quy chuẩn kỹ thuật riêng về tần số và công suất\\
  Antenna gain & Độ lợi anten & Khả năng tập trung năng lượng bức xạ vô tuyến theo hướng nhất định\\
  Transmit power & Công suất phát & Mức công suất sóng vô tuyến phát ra tại cổng anten, tính bằng dBm hoặc mW\\
  EIRP & Công suất bức xạ đẳng hướng tương đương & Công suất phát thực tế gồm công suất máy phát cộng độ lợi anten trừ suy hao cáp\\
  RSSI & Cường độ tín hiệu thu & Chỉ số biểu thị mức năng lượng tín hiệu vô tuyến đo được tại máy thu\\
  Signal-to-Noise Ratio (SNR) & Tỷ số tín hiệu trên nhiễu & Tỷ lệ giữa công suất tín hiệu hữu ích thu được so với công suất tạp âm nền\\
  Interference & Nhiễu & Tín hiệu vô tuyến ngoài ý muốn gây suy giảm chất lượng thu nhận tín hiệu\\
  Path loss & Suy hao đường truyền & Sự suy giảm mật độ công suất sóng điện từ khi truyền qua không gian\\
  Multipath & Đa đường & Hiện tượng tín hiệu truyền tới bên thu qua nhiều đường phản xạ khác nhau\\
  Fading & Pha-đinh & Sự biến thiên nhanh chóng của biên độ tín hiệu do các hiệu ứng truyền lan\\
  Modulation & Điều chế & Kỹ thuật biến đổi các đặc tính sóng mang để truyền tải dữ liệu số\\
  Coding rate & Tốc độ mã hóa & Tỷ số giữa số bit dữ liệu hữu ích thực tế trên tổng số bit truyền đi\\
  MCS & Sơ đồ điều chế và mã hóa & Bộ chỉ số xác định mức điều chế số kết hợp với tỷ lệ mã hóa kênh\\
  OFDM & Ghép kênh phân chia theo tần số trực giao & Phương pháp truyền dữ liệu trên nhiều sóng mang con trực giao nhau trong một kênh\\
  OFDMA & Đa truy cập phân chia theo tần số trực giao & Kỹ thuật phân chia các sóng mang con cho nhiều người dùng truyền nhận đồng thời\\
  MIMO & Đa đầu vào đa đầu ra & Kỹ thuật sử dụng đồng thời nhiều anten ở cả đầu phát và đầu thu vô tuyến\\
  Beamforming & Tạo búp sóng & Kỹ thuật xử lý tín hiệu nhằm tập trung búp sóng bức xạ về phía trạm nhận\\
  Throughput & Thông lượng & Tốc độ truyền dữ liệu hữu ích thực tế đến đích thành công trên mỗi giây\\
  Latency & Độ trễ & Khoảng thời gian truyền một gói dữ liệu từ điểm nguồn tới điểm đích\\

  \catrow{Điều khiển truy cập môi trường và chất lượng dịch vụ}
  Medium Access Control (MAC) & Điều khiển truy cập môi trường truyền & Lớp con điều phối quyền sử dụng môi trường kênh truyền vô tuyến dùng chung\\
  CSMA/CA & Đa truy cập tránh xung đột & Giao thức nghe sóng mang trước khi truyền kết hợp thời gian lùi ngẫu nhiên\\
  Carrier sense & Cảm nhận sóng mang & Kỹ thuật phát hiện trạng thái bận hay rỗi của môi trường truyền vô tuyến\\
  Backoff & Thời gian lùi & Thời gian tạm dừng ngẫu nhiên trước khi thử phát lại khung tin để tránh xung đột\\
  Contention Window (CW) & Cửa sổ tranh chấp & Phạm vi giá trị số nguyên dùng để lựa chọn ngẫu nhiên khoảng thời gian lùi\\
  Request to Send / Clear to Send (RTS/CTS) & Yêu cầu gửi / cho phép gửi & Thủ tục bắt tay trước nhằm bảo lưu môi trường truyền vô tuyến trước khi gửi dữ liệu\\
  Network Allocation Vector (NAV) & Vectơ phân bổ mạng & Bộ đếm thời gian được duy trì tại trạm nhằm cảm nhận sóng mang ảo\\
  Frame & Khung & Đơn vị dữ liệu giao thức ở lớp MAC được đóng gói và truyền trên sóng vô tuyến\\
  Beacon & Khung báo hiệu & Khung quản lý quảng bá định kỳ để thông báo sự hiện diện và tham số của BSS\\
  Authentication & Xác thực & Quá trình kiểm tra và xác nhận danh tính của trạm trước khi gia nhập mạng\\
  Acknowledgement (ACK) & Xác nhận & Khung điều khiển phản hồi xác nhận khung dữ liệu đã được thu nhận thành công\\
  Frame aggregation & Gộp khung & Cơ chế ghép nhiều khung dữ liệu nhỏ vào một lần truyền vô tuyến để tăng hiệu suất\\
  Quality of Service (QoS) & Chất lượng dịch vụ & Cơ chế phân loại và ưu tiên tài nguyên mạng cho các luồng dữ liệu quan trọng\\
  Access Category (AC) & Loại truy cập & Hạng mục phân loại lưu lượng theo thứ tự ưu tiên (thoại, video, thông thường, nền)\\
  Power save mode & Chế độ tiết kiệm năng lượng & Cơ chế cho phép vô tuyến chuyển sang chế độ ngủ để giảm mức tiêu thụ điện\\

  \catrow{Bảo mật}
  Confidentiality & Tính bí mật & Bảo đảm dữ liệu chỉ có thể đọc được bởi các bên được ủy quyền hợp pháp\\
  Integrity & Tính toàn vẹn & Bảo đảm dữ liệu không bị sửa đổi hay chèn thêm bất hợp pháp trong quá trình truyền\\
  Encryption & Mã hóa & Biến đổi bản rõ thành bản mã nhằm ngăn chặn việc đọc trộm nội dung thông tin\\
  Wired Equivalent Privacy (WEP) & Bảo mật tương đương mạng có dây & Chuẩn bảo mật không dây ban đầu dựa trên RC4, hiện đã lỗi thời và không an toàn\\
  Wi-Fi Protected Access (WPA) & Truy cập Wi-Fi được bảo vệ & Bộ tiêu chuẩn và chứng nhận bảo mật không dây thay thế cho WEP\\
  Advanced Encryption Standard (AES) & Chuẩn mã hóa nâng cao & Thuật toán mã hóa khối đối xứng được áp dụng bắt buộc trong WPA2 và WPA3\\
  Pre-Shared Key (PSK) & Khóa chia sẻ trước & Phương thức xác thực dùng chung mật khẩu bí mật cấu hình sẵn trên các thiết bị\\
  Simultaneous Authentication of Equals (SAE) & Xác thực đồng đẳng & Giao thức trao đổi khóa dựa trên mật khẩu có tính bảo mật chuyển tiếp trong WPA3\\
  Protected Management Frames (PMF) & Khung quản lý được bảo vệ & Cơ chế mã hóa và bảo vệ tính toàn vẹn cho các khung quản lý unicast và broadcast\\
  Pairwise Master Key (PMK) & Khóa chủ theo cặp & Khóa bí mật gốc làm cơ sở để dẫn xuất các khóa phiên tạm thời\\
  Pairwise Transient Key (PTK) & Khóa tạm thời theo cặp & Tập hợp các khóa phiên đơn hướng dùng để mã hóa và kiểm tra toàn vẹn dữ liệu\\
  Extensible Authentication Protocol (EAP) & Giao thức xác thực mở rộng & Khung giao thức hỗ trợ đóng gói nhiều phương thức xác thực người dùng khác nhau\\
  RADIUS & Dịch vụ xác thực người dùng từ xa & Giao thức AAA tiêu chuẩn dùng kết hợp với 802.1X để xác thực người dùng tập trung\\
  Virtual Private Network (VPN) & Mạng riêng ảo & Đường hầm mã hóa bảo mật truyền tải lưu lượng qua môi trường mạng không tin cậy\\
  Rogue AP & AP trái phép & Điểm truy cập trái phép được gắn vào mạng nội bộ mà không có sự phê duyệt\\
  Evil twin & AP song sinh độc hại & Điểm truy cập mạo danh phát cùng SSID của mạng hợp pháp nhằm đánh cắp thông tin\\
  Wireless Intrusion Prevention System (WIPS) & Hệ thống ngăn chặn xâm nhập không dây & Hệ thống giám sát phổ tần vô tuyến để phát hiện và ngăn chặn các cuộc tấn công\\

  \catrow{Triển khai và vận hành}
  Site survey & Khảo sát hiện trường & Hoạt động đo đạc thực địa và mô phỏng vùng phủ sóng trước khi triển khai mạng\\
  Heat map & Bản đồ nhiệt & Bản đồ biểu diễn trực quan cường độ tín hiệu hoặc thông lượng trên mặt bằng khu vực\\
  Coverage hole & Vùng lõm sóng & Khu vực có cường độ tín hiệu vô tuyến yếu dưới ngưỡng yêu cầu hoạt động\\
  Capacity & Dung lượng & Khả năng xử lý tổng lưu lượng và đáp ứng số lượng kết nối đồng thời của hệ thống\\
  Cell overlap & Độ chồng lấn ô & Phần không gian phủ sóng giao thoa giữa các AP lân cận để hỗ trợ chuyển vùng êm dịu\\
  Channel plan & Quy hoạch kênh & Kế hoạch phân bổ tần số kênh truyền cho các AP nhằm triệt tiêu nhiễu đồng kênh\\
  Airtime & Thời gian chiếm kênh & Khoảng thời gian mà một trạm hoặc AP chiếm dụng môi trường truyền vô tuyến\\
  Channel utilization & Mức sử dụng kênh & Tỷ lệ phần trăm thời gian mà kênh vô tuyến ở trạng thái bận do truyền nhận tín hiệu\\
  Retry rate & Tỷ lệ truyền lại & Tỷ lệ phần trăm các khung dữ liệu phải gửi lại do xảy ra lỗi hoặc xung đột\\
  Packet capture & Bắt gói tin & Kỹ thuật thu thập và ghi lại toàn bộ khung dữ liệu vô tuyến để chẩn đoán lỗi\\
  Monitor mode & Chế độ giám sát & Chế độ hoạt động cho phép card mạng thu nhận mọi khung dữ liệu trên kênh mà không cần liên kết\\
  Firmware & Phần mềm điều khiển & Phần mềm nhúng chuyên dụng điều khiển hoạt động phần cứng của thiết bị mạng\\
\end{xltabular}
\endgroup
